problem:  
Let \( (M^n, g) \) be a compact, simply connected Riemannian manifold with **quasipositive curvature** (\(\sec \ge 0\) everywhere, \(\sec > 0\) somewhere). Suppose \( M \) admits an **effective isometric action** by a compact connected Lie group \(G\) such that the action has **cohomogeneity two**, i.e. the orbit space \(M/G\) is a two–dimensional Alexandrov space of curvature \(\ge 0\).  In the **positively curved** case, the work of Grove–Wilking–Ziller and others shows that only finitely many **weighted orbit space configurations** (polygonal disks with labeled edges encoding isotropy types and normal slice data) can occur.  

**Question:**  
In the **quasipositive curvature** case, does there exist an effective cohomogeneity-two action on a simply connected manifold whose **weighted orbit space** \( (M/G,\, \text{labels}) \) is *not* equivalent (up to combinatorial equivalence of labeled polygons) to any of those known positively curved configurations?  
Equivalently, can local degeneracies of curvature (regions of zero sectional curvature) give rise to **new isotropy label patterns**—that is, new sequences of slice–representation weights along the boundary of \(M/G\)—which are not realizable in any positively curved example?

Why is it a "good" problem:  
- **Profound insight:**  
  The problem isolates the subtle boundary between strictly positive and merely quasipositive curvature in the presence of symmetry. While positive curvature enforces strong combinatorial restrictions on isotropy and face structure of orbit spaces, it is entirely unknown whether one can relax positivity slightly and obtain new isotropy configurations. Answering this requires understanding the transmission of local curvature positivity into global convexity and isotropy constraints on the orbit space—a delicate and largely unexplored phenomenon.
  
- **Bridge between fields:**  
  It demands an interplay between **Riemannian geometry** (sectional curvature and convexity), **Alexandrov geometry** (structure of orbit spaces with curvature bounded below), and **equivariant topology** (combinatorial classification of isotropy data). Progress would likely involve developing new compatibility conditions between curvature bounds and the combinatorial geometry of group actions, uniting techniques from these three areas.
  
- **Extreme simplicity and conceptual depth:**  
  The statement is simple: can quasipositive curvature generate *new orbit-type graphs*? Yet its resolution would illuminate how minimal positivity influences global orbit–type geometry, pressing at the boundary of known classification theory.
  
- **Potential impact:**  
  A positive example would uncover completely new manifolds with quasipositive curvature and high symmetry, overturning current intuition from positive curvature rigidity. A negative result would generalize the known rigidity of orbit structures beyond the strictly positive domain. Either outcome would substantially deepen our understanding of geometric–topological constraints in intermediate curvature regimes.  

Thus, this problem precisely captures a **sharp and open rigidity question** about how curvature positivity interacts with isotropy data in cohomogeneity‑two manifolds—conceptually simple but demanding deep methods from differential geometry, metric geometry, and transformation group theory—making it a genuinely **good research problem**.